\chapter{Fungsi}\label{ch:g10:functions}

\omterm{def:g10:functions:function}{Fungsi} adalah sebuah mesin: sebuah bilangan masuk, sebuah bilangan keluar,
dan masukan yang sama selalu memberi keluaran yang sama. \omterm{def:g10:functions:function}{Fungsi} merupakan
objek utama analisis yang akan dibangun pada tahun-tahun mendatang
(mulai \cref{ch:g11:quad}); bab ini menyiapkan kosakatanya ---
\omterm{def:g10:functions:function}{daerah asal}, \omterm{def:g10:functions:function}{peta}, \omterm{def:g10:functions:graph}{grafik}, variasi --- dan melatih keterampilan pokoknya,
yaitu membaca keterangan dari sebuah \omterm{def:g10:functions:graph}{grafik}.

\section{Kosakata: peta dan prapeta}

\begin{definition}[Fungsi]\label{def:g10:functions:function}
Misalkan $D$ sebuah himpunan \omterm{def:g10:numbers:sets}{bilangan real}. Sebuah \emph{fungsi}\index{fungsi}
$f$ yang terdefinisi pada $D$ mengaitkan setiap bilangan $x \in D$ dengan
tepat satu \omterm{def:g10:numbers:sets}{bilangan real}, ditulis $f(x)$ dan disebut \emph{peta}\index{peta}
dari $x$. Himpunan $D$ adalah \emph{daerah asal}\index{daerah asal} $f$.
Kita tulis
\[
f : x \longmapsto f(x).
\]
Jika $f(x) = y$, maka $x$ adalah \emph{prapeta}\index{prapeta} dari $y$.
\end{definition}

\begin{remark}
Setiap $x$ di \omterm{def:g10:functions:function}{daerah asal} mempunyai tepat satu \omterm{def:g10:functions:function}{peta}, tetapi sebuah bilangan
$y$ bisa mempunyai beberapa \omterm{def:g10:functions:function}{prapeta}, atau tidak sama sekali. Untuk
$f(x) = x^2$: \omterm{def:g10:functions:function}{peta} dari $3$ adalah $9$, dan
$9$ mempunyai dua \omterm{def:g10:functions:function}{prapeta}, yaitu $3$ dan $-3$; bilangan $-4$ tidak berprapeta.
\end{remark}

\begin{example}[Mencari daerah asal]\label{ex:g10:functions:domain}
Bila sebuah \omterm{def:g10:functions:function}{fungsi} diberikan oleh rumus, \omterm{def:g10:functions:function}{daerah asalnya} adalah himpunan $x$
yang membuat rumus itu bermakna.
\begin{itemize}
  \item $f(x) = 3x^2 - 5x + 1$ bermakna untuk setiap $x$: \omterm{def:g10:functions:function}{daerah asalnya}
        $\R$.
  \item $g(x) = \dfrac{1}{x - 2}$ menuntut $x - 2 \neq 0$: \omterm{def:g10:functions:function}{daerah asalnya}
        semua \omterm{def:g10:numbers:sets}{bilangan real} kecuali $2$.
  \item $h(x) = \sqrt{x - 3}$ menuntut $x - 3 \geq 0$: \omterm{def:g10:functions:function}{daerah asalnya}
        $\intco{3}{+\infty}$.
\end{itemize}
\end{example}

\begin{example}[Menghitung peta dan prapeta]\label{ex:g10:functions:images}
Misalkan $f(x) = x^2 - 4x + 3$, terdefinisi pada $\R$.

\emph{\omterm{def:g10:functions:function}{Peta} dari $5$:} substitusikan $x = 5$:
$f(5) = 25 - 20 + 3 = 8$.

\emph{\omterm{def:g10:functions:function}{Prapeta} dari $3$:} selesaikan $f(x) = 3$:
\[
x^2 - 4x + 3 = 3
\quad\Longleftrightarrow\quad
x^2 - 4x = 0
\quad\Longleftrightarrow\quad
x(x - 4) = 0,
\]
jadi \omterm{def:g10:functions:function}{prapeta} dari $3$ adalah $0$ dan $4$.
\end{example}

\section{Grafik sebuah fungsi}

\begin{definition}[Grafik]\label{def:g10:functions:graph}
Pada sebuah sistem koordinat, \emph{grafik}\index{grafik} (atau \emph{kurva})
$f$ adalah himpunan semua titik $(x, f(x))$ untuk $x$ di \omterm{def:g10:functions:function}{daerah asalnya}.
Dengan kata lain, sebuah titik $(x, y)$ terletak pada grafik tepat ketika
$y = f(x)$.
\end{definition}

\begin{omfigure}
\begin{tikzpicture}
\begin{axis}[omaxis,
  width=8.6cm, height=6.0cm,
  xmin=-2.6, xmax=4.6, ymin=-3.4, ymax=4.4,
  xtick={-2,-1,1,2,3,4}, ytick={-3,-2,-1,1,2,3,4}]
  \addplot[omDef, thick, domain=-2.2:4.2, samples=120]
    {0.25*x^3 - 0.75*x^2 - x + 2};
  \draw[dashed, omThm] (axis cs:3,0) -- (axis cs:3,-1)
    -- (axis cs:0,-1);
  \fill[omThm] (axis cs:3,-1) circle (1.5pt);
  \node[omThm, font=\small, below right] at (axis cs:3,-1.05)
    {$(3,\,f(3))$};
\end{axis}
\end{tikzpicture}

{\small Membaca sebuah \omterm{def:g10:functions:function}{peta} pada \omterm{def:g10:functions:graph}{grafik}: mulailah dari $x = 3$ pada sumbu
mendatar, naiklah tegak lurus sampai kurva, lalu bergeraklah mendatar ke
sumbu tegak untuk membaca $f(3) = -1$.}
\end{omfigure}

\begin{method}[Membaca grafik]\label{met:g10:functions:reading}
Pada \omterm{def:g10:functions:graph}{grafik} $f$:
\begin{enumerate}
  \item \emph{\omterm{def:g10:functions:function}{peta} dari $a$}: bergeraklah tegak lurus dari $x = a$ pada sumbu
        mendatar sampai kurva, lalu mendatar ke sumbu tegak; nilai yang
        terbaca di sana adalah $f(a)$;
  \item \emph{\omterm{def:g10:functions:function}{prapeta} dari $b$}: tariklah garis mendatar $y = b$; \omterm{def:g10:functions:function}{prapetanya}
        adalah koordinat $x$ dari semua titik potong garis itu
        dengan kurva;
  \item \emph{penyelesaian $f(x) = k$}: sama dengan \omterm{def:g10:functions:function}{prapeta} dari $k$;
  \item \emph{penyelesaian $f(x) \leq k$}: nilai $x$ yang membuat kurva
        berada pada atau di bawah garis $y = k$.
\end{enumerate}
\end{method}

\begin{omfigure}
\begin{tikzpicture}
\begin{axis}[omaxis,
  width=8.6cm, height=6.0cm,
  xmin=-3.4, xmax=3.4, ymin=-2.4, ymax=4.4,
  xtick={-2,2}, xticklabels={$x_1$,$x_2$}, ytick={2}, yticklabels={$k$}]
  \addplot[omDef, thick, domain=-3:3, samples=120] {0.5*x^2};
  \addplot[omThm, thick, domain=-3.2:3.2] {2};
  \fill (axis cs:-2,2) circle (1.5pt);
  \fill (axis cs:2,2) circle (1.5pt);
  \draw[dashed] (axis cs:-2,2) -- (axis cs:-2,0);
  \draw[dashed] (axis cs:2,2) -- (axis cs:2,0);
\end{axis}
\end{tikzpicture}

{\small Menyelesaikan $f(x) = k$ secara grafis: penyelesaian $x_1$ dan $x_2$
adalah absis titik potong kurva dengan garis mendatar
$y = k$. Di sini $f(x) \leq k$ berlaku pada $\intcc{x_1}{x_2}$, yaitu tempat
kurva berada di bawah garis itu.}
\end{omfigure}

\begin{example}\label{ex:g10:functions:membership}
Apakah titik $A(2, 5)$ terletak pada \omterm{def:g10:functions:graph}{grafik} $f(x) = x^2 + 1$? Hitunglah
$f(2) = 4 + 1 = 5$: ya, sebab $f(2)$ sama dengan koordinat $y$ titik $A$.
Titik $B(3, 8)$ tidak terletak pada \omterm{def:g10:functions:graph}{grafik} itu, sebab $f(3) = 10 \neq 8$.
\end{example}

\section{Variasi sebuah fungsi}

\begin{definition}[Naik, turun]\label{def:g10:functions:variations}
Misalkan $f$ terdefinisi pada sebuah \omterm{def:g10:numbers:interval}{interval} $I$.
\begin{itemize}
  \item $f$ disebut \emph{naik}\index{fungsi naik} pada $I$ apabila
        untuk semua $u, v \in I$, jika $u < v$ maka $f(u) \leq f(v)$:
        keluarannya membesar seiring masukannya, dan \omterm{def:g10:functions:graph}{grafiknya} mendaki dari
        kiri ke kanan.
  \item $f$ disebut \emph{turun}\index{fungsi turun} pada $I$ apabila
        $u < v$ mengakibatkan $f(u) \geq f(v)$: \omterm{def:g10:functions:graph}{grafiknya} menurun dari kiri
        ke kanan.
\end{itemize}
Dengan ketaksamaan tegas $f(u) < f(v)$ (masing-masing $>$) kita katakan
\emph{tegas} naik (masing-masing turun).
\end{definition}

\begin{definition}[Tabel variasi]\label{def:g10:functions:table}
Sebuah \emph{tabel variasi}\index{tabel variasi} merangkum pada \omterm{def:g10:numbers:interval}{interval} mana
saja $f$ \omterm{def:g10:functions:variations}{naik} dan \omterm{def:g10:functions:variations}{turun}, dengan panah $\nearrow$ dan
$\searrow$, serta mencatat nilai $f$ pada titik-titik baliknya.
\end{definition}

\begin{example}\label{ex:g10:functions:table}
Tinjaulah \omterm{def:g10:functions:function}{fungsi} yang \omterm{def:g10:functions:graph}{grafiknya} digambar di bawah ini pada $\intcc{-2}{4}$.

\begin{omfigure}
\begin{tikzpicture}
\begin{axis}[omaxis,
  width=8.6cm, height=5.4cm,
  xmin=-2.8, xmax=4.6, ymin=-2.4, ymax=3.8,
  xtick={-2,1,4}, ytick={-1.5,3}, yticklabels={$-1.5$,$3$}]
  \addplot[omDef, thick, domain=-2:4, samples=120] {-0.5*(x-1)^2 + 3};
  \fill (axis cs:-2,-1.5) circle (1.5pt);
  \fill (axis cs:4,-1.5) circle (1.5pt);
  \fill (axis cs:1,3) circle (1.5pt);
\end{axis}
\end{tikzpicture}

{\small Sebuah \omterm{def:g10:functions:function}{fungsi} pada $\intcc{-2}{4}$, mula-mula \omterm{def:g10:functions:variations}{naik} lalu \omterm{def:g10:functions:variations}{turun}.}
\end{omfigure}

Dari \omterm{def:g10:functions:graph}{grafiknya}: $f$ \omterm{def:g10:functions:variations}{naik} dari $f(-2) = -1.5$ sampai \omterm{def:g10:functions:extrema}{maksimumnya}
$f(1) = 3$, lalu \omterm{def:g10:functions:variations}{turun} sampai $f(4) = -1.5$. \omterm{def:g10:functions:table}{Tabel variasinya} adalah
\begin{center}
\renewcommand{\arraystretch}{1.3}
\begin{tabular}{c|ccccc}
$x$ & $-2$ & & $1$ & & $4$ \\
\hline
$f$ & $-1.5$ & $\nearrow$ & $3$ & $\searrow$ & $-1.5$ \\
\end{tabular}
\end{center}
\end{example}

\begin{example}[Membuktikan sebuah variasi]\label{ex:g10:functions:proof}
Tunjukkan bahwa $f(x) = 3x + 1$ tegas \omterm{def:g10:functions:variations}{naik} pada $\R$. Ambil sebarang
$u < v$ lalu bandingkan \omterm{def:g10:functions:function}{petanya}:
\[
f(v) - f(u) = (3v + 1) - (3u + 1) = 3(v - u) > 0,
\]
sebab $v - u > 0$. Jadi $f(u) < f(v)$: \omterm{def:g10:functions:function}{fungsi} itu tegas
\omterm{def:g10:functions:variations}{naik}. Perhitungan yang sama dengan gradien negatif, misalnya
$g(x) = -2x + 5$, memberi $g(v) - g(u) = -2(v-u) < 0$: tegas \omterm{def:g10:functions:variations}{turun}.
\end{example}

\section{Nilai ekstrem}

\begin{definition}[Maksimum, minimum]\label{def:g10:functions:extrema}
Misalkan $f$ terdefinisi pada $D$ dan $a \in D$. Nilai $f(a)$ disebut
\emph{maksimum}\index{maksimum} $f$ pada $D$ apabila $f(x) \leq f(a)$ untuk
semua $x \in D$; nilai itu disebut \emph{minimum}\index{minimum} apabila
$f(x) \geq f(a)$ untuk semua $x \in D$. Pada \omterm{def:g10:functions:graph}{grafik}, keduanya adalah titik
tertinggi dan titik terendah kurvanya.
\end{definition}

\begin{example}\label{ex:g10:functions:extremum}
Tunjukkan bahwa $f(x) = (x - 1)^2 + 2$ mempunyai \omterm{def:g10:functions:extrema}{minimum} $2$ pada $\R$, yang
dicapai di $x = 1$. Untuk setiap $x$, kuadrat $(x-1)^2$ bernilai $\geq 0$,
jadi $f(x) = (x-1)^2 + 2 \geq 2$; dan $f(1) = 0 + 2 = 2$, sehingga nilai $2$
benar-benar tercapai. Kedua fakta itu bersama membuktikan bahwa $2$ \omterm{def:g10:functions:extrema}{minimumnya}.
\end{example}

\begin{omfigure}
\begin{tikzpicture}
\begin{axis}[omaxis,
  width=7.6cm, height=5.2cm,
  xmin=-1.8, xmax=3.8, ymin=-0.9, ymax=6.4,
  xtick={1}, ytick={2}]
  \addplot[omDef, thick, domain=-1.0:3.0, samples=100] {(x-1)^2 + 2};
  \addplot[omThm, dashed, domain=-1.6:3.6] {2};
  \fill (axis cs:1,2) circle (1.5pt);
  \node[below right, font=\small] at (axis cs:1,2) {$(1,2)$};
\end{axis}
\end{tikzpicture}

{\small Kurva $f(x) = (x-1)^2 + 2$ tidak pernah \omterm{def:g10:functions:variations}{turun} di bawah garis putus-putus
$y = 2$, dan menyentuhnya di $x = 1$: \omterm{def:g10:functions:extrema}{minimum} $f$ adalah $2$.}
\end{omfigure}

\section{Latihan}

\begin{exercise}[$\star$]\label{exo:g10:functions:1}
Misalkan $f(x) = 2x^2 - 3x + 1$. Hitunglah \omterm{def:g10:functions:function}{peta} dari $0$, $2$, $-1$ dan
$\frac12$.
\end{exercise}

\begin{exercise}[$\star$]\label{exo:g10:functions:2}
Berikan \omterm{def:g10:functions:function}{daerah asal} setiap \omterm{def:g10:functions:function}{fungsi} berikut:
\[
f(x) = 5x - 2, \qquad
g(x) = \frac{3}{x + 4}, \qquad
h(x) = \sqrt{2x - 6}, \qquad
k(x) = \frac{1}{x^2 + 1}.
\]
\end{exercise}

\begin{exercise}[$\star$]\label{exo:g10:functions:3}
Misalkan $f(x) = x^2 - 2x$. Carilah semua \omterm{def:g10:functions:function}{prapeta} dari $0$, dari $3$, dan dari $-1$.
\end{exercise}

\begin{exercise}[$\star$]\label{exo:g10:functions:4}
Apakah titik $A(-1, 4)$ terletak pada \omterm{def:g10:functions:graph}{grafik} $f(x) = x^2 - 2x + 1$?
Dan titik $B(2, 1)$? Berikan alasan lewat perhitungan.
\end{exercise}

\begin{exercise}[$\star$]\label{exo:g10:functions:5}
Sebuah \omterm{def:g10:functions:function}{fungsi} $g$ pada $\intcc{-3}{5}$ mempunyai \omterm{def:g10:functions:table}{tabel variasi}
\begin{center}
\renewcommand{\arraystretch}{1.3}
\begin{tabular}{c|ccccccc}
$x$ & $-3$ & & $0$ & & $3$ & & $5$ \\
\hline
$g$ & $1$ & $\searrow$ & $-2$ & $\nearrow$ & $4$ & $\searrow$ & $0$ \\
\end{tabular}
\end{center}
\begin{enumerate}
  \item Berapakah \omterm{def:g10:functions:extrema}{maksimum} dan \omterm{def:g10:functions:extrema}{minimum} $g$ pada $\intcc{-3}{5}$?
  \item Bandingkan $g(-1)$ dan $g(-0.5)$ tanpa mengetahui nilainya.
  \item Paling banyak berapa penyelesaian yang dipunyai \omterm{def:g10:algebra:equation}{persamaan} $g(x) = 0$
        pada setiap \omterm{def:g10:numbers:interval}{interval} dalam tabel itu?
\end{enumerate}
\end{exercise}

\begin{exercise}[$\star\star$]\label{exo:g10:functions:6}
Misalkan $f(x) = -2x + 7$. Tunjukkan, dengan membandingkan $f(u)$ dan $f(v)$
untuk $u < v$, bahwa $f$ tegas \omterm{def:g10:functions:variations}{turun} pada $\R$.
\end{exercise}

\begin{exercise}[$\star\star$]\label{exo:g10:functions:7}
Tunjukkan bahwa \omterm{def:g10:functions:function}{fungsi} $f(x) = x^2 + 6x + 11$ mempunyai \omterm{def:g10:functions:extrema}{minimum} pada $\R$,
lalu berikan nilainya dan tempat nilai itu dicapai. (Petunjuk: tulislah
$f(x) = (x + 3)^2 + c$ dengan konstanta $c$ yang tepat.)
\end{exercise}

\begin{exercise}[$\star\star$]\label{exo:g10:functions:8}
Sebuah kebun berbentuk persegi panjang berkeliling $40$ m. Misalkan $x$
lebarnya, dalam meter.
\begin{enumerate}
  \item Nyatakan panjangnya, lalu luasnya $A(x)$, sebagai \omterm{def:g10:functions:function}{fungsi} dari $x$.
        Untuk $x$ yang mana hal ini bermakna?
  \item Hitunglah $A(4)$, $A(8)$, $A(10)$, $A(12)$ dan $A(16)$. Apa dugaanmu
        tentang bentuk yang luasnya terbesar?
\end{enumerate}
\end{exercise}

\begin{exercise}[$\star\star$]\label{exo:g10:functions:9}
Misalkan $f(x) = \dfrac{1}{x^2 + 1}$, terdefinisi pada $\R$.
\begin{enumerate}
  \item Tunjukkan bahwa $0 < f(x) \leq 1$ untuk setiap \omterm{def:g10:numbers:sets}{bilangan real} $x$.
  \item Apakah $f$ mempunyai \omterm{def:g10:functions:extrema}{maksimum} pada $\R$? \omterm{def:g10:functions:extrema}{Minimum}? Berikan alasannya.
\end{enumerate}
\end{exercise}

\begin{exercise}[$\star\star\star$]\label{exo:g10:functions:10}
Misalkan $f(x) = x^2$ pada $\R$.
\begin{enumerate}
  \item Misalkan $0 \leq u < v$. Faktorkan $f(v) - f(u)$ lalu simpulkan bahwa
        $f$ tegas \omterm{def:g10:functions:variations}{naik} pada $\intco{0}{+\infty}$.
  \item Sesuaikan alasannya untuk menunjukkan bahwa $f$ tegas \omterm{def:g10:functions:variations}{turun} pada
        $\intoc{-\infty}{0}$.
\end{enumerate}
\end{exercise}

\section{Soal: Mesin yang disuapi keluarannya sendiri}

\begin{problem}[{Soal akhir pekan --- titik tetap, iterasi, dan
tiga mesin termasyhur: penyempurna Heron, pengkuadrat, dan mesin hujan es
yang belum terpecahkan}]
\label{pb:g10:functions:1}
\omterm{def:g10:functions:function}{Fungsi} adalah sebuah mesin: sebuah bilangan masuk, sebuah bilangan keluar.
Percobaan yang paling menarik menyuapi mesin itu dengan \emph{keluarannya
sendiri}, berulang-ulang --- lalu dua pertanyaan menjadi utama:
apakah prosesnya mengendap di suatu tempat (sebuah \emph{\omterm{pb:g10:functions:1}{titik tetap}}), dan
bagaimana variasi mesin itu menuntunnya ke sana? Salah satu mesin dalam
soal ini sudah dua ribu tahun mengasah akar kuadrat; mesin yang lain menjaga
soal tak terpecahkan yang paling termasyhur dalam matematika dasar.

\medskip
\noindent\textbf{Bagian I --- Mesin Heron.}
Tinjaulah \omterm{def:g10:functions:function}{fungsi}
\[
f(x) = \frac12\left(x + \frac2x\right).
\]
\begin{enumerate}
  \item Berikan \omterm{def:g10:functions:function}{daerah asal} $f$, lalu hitunglah $f(1)$, $f(2)$ dan
        $f\!\left(\frac32\right)$ sebagai pecahan eksak.
  \item Carilah semua \omterm{def:g10:functions:function}{prapeta} dari $\frac32$: selesaikan
        $f(x) = \frac32$ (hilangkan penyebutnya lalu pakai
        metode \cref{pb:g10:algebra:1}).
  \item Sebuah \emph{titik tetap}\index{titik tetap} $f$ adalah
        bilangan yang tidak diubah oleh mesin itu: $f(x) = x$. Carilah kedua
        \omterm{pb:g10:functions:1}{titik tetap} $f$. Kenalan lama yang mana menjaga
        \omterm{pb:g10:functions:1}{titik tetap} positifnya?
  \item Mulai dari $x = 1$, suapkan keluaran mesin itu kembali ke dalamnya
        tiga kali lalu catat hasil eksaknya. Barisan mana pada
        soal akhir pekan tentang sifat \omterm{ex:g10:numbers:classify}{irasional} di jilid sebelumnya yang baru
        saja dibangun kembali --- dan apa itu ``resep Heron'' dalam bahasa \omterm{def:g10:functions:function}{fungsi}?
  \item Pada satu gambar, sketsalah \omterm{def:g10:functions:graph}{grafik} $f$ (untuk $x > 0$)
        dan garis diagonal $y = x$. Di manakah letak \omterm{pb:g10:functions:1}{titik
        tetapnya} pada gambar itu? (Bandingkan dengan
        soal akhir pekan tentang dua termometer di jilid sebelumnya: suhu yang
        terbaca sama pada kedua skala adalah gagasan yang sama.)
\end{enumerate}

\medskip
\noindent\textbf{Bagian II --- Mengapa mesin itu tak mungkin meleset.}
\begin{enumerate}[resume]
  \item Untuk $0 < u < v$, tunjukkan dengan menyamakan
        penyebutnya bahwa
        \[
        f(v) - f(u) = \frac{(v - u)(uv - 2)}{2uv} ,
        \]
        lalu simpulkan variasi $f$ pada
        $\intoo{0}{+\infty}$: \omterm{def:g10:functions:variations}{turun} sampai suatu titik, lalu
        \omterm{def:g10:functions:variations}{naik} --- titik yang mana? (Teknik yang sama seperti pada
        \cref{exo:g10:functions:10}.)
  \item Simpulkan bahwa pada $\intoo{0}{+\infty}$ \omterm{def:g10:functions:function}{fungsi} $f$ mempunyai
        \omterm{def:g10:functions:extrema}{minimum}, yang nilainya tepat $\sqrt2$, dicapai di
        $\sqrt2$. Apa artinya hal ini bagi \emph{setiap} keluaran
        mesin itu (untuk masukan positif)?
  \item Tunjukkan bahwa jika $x > \sqrt2$ maka
        $\sqrt2 < f(x) < x$: terkaan yang kebesaran selalu diperbaiki,
        tidak pernah kelewatan. (Untuk $f(x) < x$, bandingkan $\frac2x$ dengan
        $x$.)
  \item Rangkailah pertanyaan 6--8 menjadi \omterm{def:g10:functions:table}{tabel variasi} $f$
        pada $\intoo{0}{+\infty}$
        (\cref{def:g10:functions:table}), lengkap dengan \omterm{def:g10:functions:extrema}{minimumnya}.
  \item Mesin kedua: $g(x) = x^2$. \omterm{pb:g10:functions:1}{Titik tetapnya} adalah
        $0$ dan $1$ (periksa). Iterasikan $g$ empat kali dari
        $x = 0.9$, lalu empat kali dari $x = 1.1$ (dengan kalkulator,
        tiga angka desimal). Satu \omterm{pb:g10:functions:1}{titik tetap} menarik, yang lain
        menolak: yang mana yang mana?
\end{enumerate}

\medskip
\noindent\textbf{Bagian III --- Mesin hujan es.}
Pada \omterm{def:g10:numbers:sets}{bilangan bulat} positif, definisikan: $h(n) = \frac n2$ jika $n$ genap,
dan $h(n) = 3n + 1$ jika $n$ ganjil. Bilangan memantul di bawah $h$ seperti
butir hujan es di dalam awan badai.
\begin{enumerate}[resume]
  \item Hitunglah seluruh penerbangan $7$ di bawah $h$, sampai ke $1$.
        Berapa langkah yang diperlukan, dan setinggi apa ia terbang?
  \item Mulailah penerbangan $27$ lalu hitunglah sepuluh langkah. (Seluruh
        penerbangannya berlangsung $111$ langkah dan memuncak di $9\,232$ ---
        padahal berangkat dari $27$.) Pelajaran apa tentang aturan
        sederhana yang diajarkan oleh $7$ dan $27$?
  \item Jelaskan mengapa setiap penerbangan yang mencapai pangkat dari $2$
        langsung terjun ke $1$, dan apa yang terjadi di $1$
        (hitunglah $h(1)$, $h(4)$, $h(2)$). Apakah $h$ mempunyai
        \omterm{pb:g10:functions:1}{titik tetap} di antara \omterm{def:g10:numbers:sets}{bilangan bulat} positif? Apa yang justru
        dibentuk oleh ujung setiap penerbangan yang teramati?
  \item \emph{Dugaan Collatz} menyatakan bahwa setiap bilangan
        awal pada akhirnya mencapai $1$. Komputer telah memeriksanya
        jauh melampaui $10^{20}$; tidak seorang pun pernah membuktikannya. Dua
        macam temuan apa yang akan \emph{menggugurkannya}? Dan mengapa
        pemeriksaan besar-besaran tetap tidak menutup perkaranya
        (soal akhir pekan tentang ramalan batang korek api di jilid sebelumnya
        pernah membunyikan lonceng ini)?
  \item ``Matematika belum matang untuk pertanyaan semacam itu,'' kata
        Paul Erd\H{o}s tentang dugaan ini. Dalam satu kalimat: apa
        yang membedakan mesin hujan es dari mesin Heron, yang
        pertanyaan 6--8 sudah menuntaskan segalanya?
\end{enumerate}

\medskip
\noindent\textbf{Bagian IV --- Buku petunjuk pemakai mesin.}
\begin{enumerate}[resume]
  \item \omterm{def:g10:functions:function}{Daerah asal} sebagai \omterm{def:g10:numbers:interval}{interval} (dalam bahasa
        \cref{pb:g10:numbers:1}): berikan \omterm{def:g10:functions:function}{daerah asal}
        $x \mapsto \sqrt{x - 3}$, \
        $x \mapsto \frac{1}{x^2 - 9}$, \
        $x \mapsto \sqrt{9 - x^2}$.
  \item Kembali ke $f$ milik Heron: carilah \omterm{def:g10:functions:function}{prapeta} eksak dari $3$
        (selesaikan $f(x) = 3$ dengan melengkapkan kuadrat sempurna).
  \item Sebuah batu dilempar ke atas: tingginya
        $H(t) = 20t - 5t^2$ meter setelah $t$ detik. Pada selang
        waktu mana rumus itu bermakna secara fisis?
        Lengkapkan kuadratnya untuk mencari tinggi \omterm{def:g10:functions:extrema}{maksimum} dan
        waktunya (\cref{def:g10:functions:extrema}), lalu berikan
        \omterm{def:g10:functions:table}{tabel variasinya}.
  \item Berapa lama batu itu berada paling sedikit $15$ m di atas tanah?
        (Selesaikan $H(t) \geq 15$ dengan \omterm{def:g10:algebra:signtable}{tabel tanda},
        \cref{met:g10:algebra:signtable}.)
  \item Penutup: soal ini menjalankan empat mesin --- penyempurna
        Heron, pengkuadrat, mesin hujan es, dan tinggi batu.
        Dalam dua atau tiga kalimat, nyatakan apa yang dimiliki setiap
        \omterm{def:g10:functions:function}{fungsi} (\omterm{def:g10:functions:function}{daerah asal}, aturan, \omterm{def:g10:functions:graph}{grafik}) dan dua
        pertanyaan besar yang diajukan soal ini kepada masing-masing (ke mana
        iterasinya mengendap? bagaimana \omterm{def:g10:functions:function}{fungsi} itu bervariasi?) ---
        pertanyaan yang akan digarap satu per satu oleh bab-bab berikutnya.

\end{enumerate}
\end{problem}
